
\subsection{Protein Crystallization}
\begin{frame}{Nucleation Barrier}
    \begin{columns}[T]
        \column{0.55\textwidth}
        \begin{itemize}
            \item \textbf{Surface Free Energy} destabilizes the nuclei $\Delta G_S = 4 \pi r^2 \gamma$
            \item \textbf{Volume (Bulk) Free Energy} stabilizes the nuclei $\Delta G_V = \frac{4}{3} \pi r^3 \Delta g_v$
            \item $\Delta g_v = \frac{\text{free energy}}{\text{unit volume}}$
            \item $\Delta G = \Delta G_S + \Delta G_V$
            \item $r^*$ = critical radius, bulk term dominates and growth is favorable
            \item the free energy maximum $\Delta G^*$ is the \textbf{nucleation barrier}
        \end{itemize}

        \column{0.45\textwidth}
        \vspace{-0.55cm}
        \begin{figure}
            \centering
            \includegraphics[width=0.9\linewidth]{../msc-thesis-code/data/dataExploration/plots/free_energy_function}
            \caption[Free Energy Nucleation Diagram]{Free Energy Nucleation Diagram}
            \label{fig:nucleationBarrier}
        \end{figure}
    \end{columns}
    \label{frame:NucleationBarrier}
\end{frame}

\begin{frame}{Vapor Diffusion}
    \begin{columns}[T]
        \column{0.55\textwidth}
\begin{itemize}
    \item Vapor diffusion changes the droplet composition over time
    \item Loss of water increases both \textbf{protein concentration} and \textbf{precipitant concentration}
    \item In the phase diagram, this means the system moves from \textbf{undersaturation} toward \textbf{supersaturation}
    \item Reaching the \textbf{nucleation zone} enables the formation of crystal nuclei
    \item Remaining in or returning to the \textbf{metastable zone} promotes growth of fewer, larger crystals
\end{itemize}
        \column{0.45\textwidth}
        \begin{figure}
            \centering
            \includegraphics[width=0.9\linewidth]{Figures/hangingDrop}
            \caption[Hanging Drop Vapor Diffusion]{Hanging Drop Vapor Diffusion}
            \label{fig:nucleationBarrier}
        \end{figure}
    \end{columns}
    \label{frame:NucleationBarrier}
\end{frame}

\begin{frame}{Compounds \& X-ray Diffraction}
    \begin{itemize}
        \item \textbf{Organic precipitants:} PEG lower protein solubility and drive the system toward supersaturation
        \item \textbf{Salts:} Salt concentration can either increase or decrease protein solubility
        \item \textbf{Buffers:} Buffers set the pH and thus shape the protein’s surface charge distribution $\rightarrow$ intermolecular interactions and crystal formation.
    \end{itemize}

    \begin{figure}
        \centering
        \includegraphics[width=0.9\linewidth]{Figures/x-ray-crystallography.pdf}
        \caption[X-ray Diffraction]{X-ray Diffraction setup}
        \label{fig:phaseDiagram}
    \end{figure}
    \label{frame:CompoundsXRay}
\end{frame}
